\newpage
\sscp{\tsc{Sumba-Havu} classification}{03-04-03}

\dirtree{%
.1 {\tsc{Bima-Lembata} *b > **β, **b}.
.2 {\tsc{Sumba-Havu} *R > **y, **R; *d/*j > **r}.
.3 {\tsc{Sumba} *z/**R/**r > \it{r}; lexical similarity \citep{Gasser2014b}}.
.4 {\tsc{Kodi-Gaura}}.
.5 {Kodi [kod]}.
.5 {Gaura [{-}{-}{-}] *mb > \it{b}; *mb > \it{b}; *ŋg > \it{g}}.
.4 {Laboya [lmy] *mb > \it{b}; *mb > \it{b}; *ŋg > \it{g}}.
.4 {\tsc{West Sumba}}.
.5 {Bukambero [{-}{-}{-}]}.
.5 {Loura [lur]}.
.5 {Tana Righu [{-}{-}{-}]}.
.5 {Weyewa [wew]}.
.5 {South Weyewa [{-}{-}{-}]}.
.5 {Waibangga [{-}{-}{-}]}.
.5 {Loli  [{-}{-}{-}]}.
.4 {\tsc{Central Sumba}}.
.5 {Mamboru [mvd] (*s = \it{s})}.
.5 {\tsc{Nuclear Central Sumba} *mb > \it{b}; *mb > \it{b}; *ŋg > \it{g}}.
.6 {Pondok [{-}{-}{-}]}.
.6 {Baliledo [{-}{-}{-}]}.
.6 {Wanukaka [wnk]}.
.6 {Anakalang [akg]}.
.4 {\tsc{East Sumba}}.
.5 {\tsc{Maderi-Praiwatana}}.
.6 {Maderi}.
.6 {Praiwatana}.
.5 {\tsc{Nuclear East Sumba} [xbr]}.
.6 {Umbu Ratu Nggai,
		Lewa, Hahar,
		Prai Paha,
		Kanatang-Kapunduk,
		Kambera,
		Tabundung-Karera,
		Melolo (a.k.a. Rindi)}.
%.6 {Umbu Ratu Nggai}.
%.6 {Lewa}.
%.6 {Hahar}.
%.6 {Prai Paha}.
%.6 {Kanatang-Kapunduk}.
%.6 {Kambera}.
%.6 {Tabundung-Karera}.
%.6 {Melolo (a.k.a. Rindi)}.
.5 {Mangili [{-}{-}{-}]}.
.3 {\tsc{Havu-Dhao} *z > ʄ /{\#}{\gap}; *h/*R/*q/*w > {\0};
										*ə > \it{e; ə}; %/{\#}σ\und{σ}{\#};
										%*u > u; o/[both syllables are *u];
										%*-ay > \it{i; e} /[harmony of V height with mid V];
										*-aw > \it{o}; *-uy > \it{i}; \\ \hp{\tsc{h}} *-iw > i;
										*i(C)a/*a(C)i > **ə(C)i; *u(C)a/*a(C)u > **ə(C)u; (*ñ = \it{ɲ})}.
.4 {\tsc{Havu-Raijua} *s > \it{h}; **β > /v/ → [v]{\tl}[w]{\tl}[β]; VSO ergative marked}.
.5 {Havu [hvn]}.
.5 {Raijua [{-}{-}{-}]}.
.4 {Dhao [nfa] **β > \it{h}; *s > \it{ʧ} /{\#}{\gap} and /V{\gap}V\tsc{[+high]}; SVO \tsc{nom-acc}}.
}